\section{Stationary Correlated-Stream Realizations}

The block experiment has a direct stationary interpretation. Suppose first
that integer durations $\ell_j\asymp j^r$ are used. Concatenate i.i.d. blocks
with modes drawn from $q$ and durations $\ell_j$. Initialize the block
containing time zero from the equilibrium size-biased distribution and choose
its age uniformly within the block. This produces a stationary semi-Markov
stream. Equation~\eqref{eq:occupancy} shows that its one-time covariance is
$\sum_j\lambda_je_je_j^\top$.

Let $N_B=\sum_{k=1}^B\ell_{J_k}$ be the raw duration of $B$ completed blocks.
Then
\begin{equation}
\E N_B=\mu B,
\qquad
N_B/B\longrightarrow\mu\quad\text{almost surely}.
\label{eq:renewal-slln}
\end{equation}
Consequently, indexing by expected raw observations changes constants but not
any exponent. If $r<a-1$, the duration has finite variance and $N_B$ has
ordinary square-root concentration. We keep Theorem~\ref{thm:exact-risk}
indexed by completed innovations because it yields an exact finite-sample
identity and avoids boundary effects from a partially observed final block.
Almost-sure equivalence by itself does not imply an expected-risk equivalence:
rare, unusually long blocks could dominate. The following result controls that
possibility directly.

\begin{theorem}[Fixed raw horizon and the inspection paradox]
\label{thm:raw-horizon}
Assume Equation~\eqref{eq:power-laws}, $\sigma=0$, integer dwell times
$\tau_j\ge1$, and
\begin{equation}
1<b<a+r+1.
\label{eq:raw-horizon-conditions}
\end{equation}
First start the concatenated stream at a renewal boundary and observe exactly
$N$ raw examples.  For every $r\ge0$, uniformly for $N,M\ge2$,
\begin{equation}
\E\risk(\widehat\theta^{(N)},\theta)
\asympconst
M^{-(b-1)}+N^{-(b-1)/(a+r)}.
\label{eq:raw-horizon-rate}
\end{equation}
No finite-variance assumption on the dwell time is needed.

Now initialize the stream from its stationary equilibrium renewal law.  If
$r>0$, then
\begin{equation}
\E\risk(\widehat\theta^{(N)}_{\rm eq},\theta)
\asympconst
M^{-(b-1)}+N^{-(b-1)/(a+r)}+N^{-(a-1)/r}.
\label{eq:stationary-raw-rate}
\end{equation}
For $r=0$, the last term is absent.  Thus equilibrium initialization preserves
the boundary-start exponent when $(b-1)/(a+r)<(a-1)/r$, but otherwise the
forward recurrence time creates a slower inspection-paradox phase.
\end{theorem}

The proof follows individual mode-hitting times rather than concentrating the
total number of renewals.  If $H_j$ is the raw start time of the first block in
mode $j$, then
\begin{equation}
\E H_j=\mu/q_j-\tau_j\le\mu/q_j,
\qquad
(1-q_j)^N\le\Prob(H_j\ge N)
\le\min\{1,\mu/(Nq_j)\}.
\label{eq:hitting-summary}
\end{equation}
Summing these inequalities against $s_j$ gives
Equation~\eqref{eq:raw-horizon-rate} for all $r$.  Under stationary
initialization, the residual duration $R$ of the initial size-biased block has
\begin{equation}
\Prob(R\ge N)\asympconst N^{-(a-1)/r},
\label{eq:inspection-tail-main}
\end{equation}
which yields the additional term in Equation~\eqref{eq:stationary-raw-rate}.
The technical supplement gives matching upper and lower bounds.

A Markov realization makes the nonseparable covariance explicit. Consider
states $(j,s)$ with $s\in\{-1,+1\}$ and transition kernel
\begin{align}
P((j,s),(k,t))
={}&\left(1-\tau_j^{-1}\right)
\1\{(j,s)=(k,t)\}
+\tau_j^{-1}\frac{q_k}{2}.
\label{eq:markov-kernel}
\end{align}

\begin{proposition}[Reversible mode-dependent persistence]
\label{prop:markov}
The chain in Equation~\eqref{eq:markov-kernel} is reversible with stationary
law $\pi_{j,s}=\lambda_j/2$. For $X_t=S_te_{J_t}$,
\begin{equation}
\E[X_tX_{t+\ell}^\top]
=\sum_{j\ge1}\lambda_j
\left(1-\tau_j^{-1}\right)^\ell e_je_j^\top.
\label{eq:mode-autocovariance}
\end{equation}
Unless all $\tau_j$ are equal, this covariance cannot be written as a scalar
time kernel multiplied by the spatial covariance.
\end{proposition}

\begin{proof}
For distinct states,
\[
\pi_{j,s}P((j,s),(k,t))
=\frac{\lambda_jq_k}{4\tau_j}
=\frac{Zq_jq_k}{4},
\]
which is symmetric in $(j,s)$ and $(k,t)$; self-transitions are immediate.
For the covariance, the only contribution correlated with $X_t$ is the event
that no refresh occurs during the next $\ell$ steps. A refresh resamples a
mean-zero sign, giving Equation~\eqref{eq:mode-autocovariance}.
\end{proof}

The proposition explains why a single sample-correlation matrix is
insufficient here: each spatial mode carries a different temporal kernel.
The exact noisy theorem uses deterministic or block-averaged durations;
extending its sharp variance constants to geometric dwell times is discussed
under limitations.

